\chapter*{РЭФЕРАТ}
\addcontentsline{toc}{chapter}{РЭФЕРАТ}
\vspace{-1cm}

Дыпломная работа: 71 старонка, 14 рысункаў, 2 табліцы, 30 крыніц, 2 дадаткі.

\textbf{Ключавыя словы:} КВАНТАВАЯ КРОПКА, МЕТАД МОЦНАЙ СУВЯЗІ, СУПЕРЭЛІПС,
НІЗКАЭНЕРГЕТЫЧНЫ СПЕКТР, СУРАГАТНАЯ АПРАКСІМАЦЫЯ,
ГРЭБЕНЕВАЯ РЭГРЭСІЯ, МНАГАСЛОЙНЫ ПЕРЦЭПТРОН,
СТРУКТУРАВАНАЯ ПРАВЕРКА.\par

\textbf{Аб'ект даследавання} --- квантавыя кропкі суперэліптычнай формы на квадратнай рашотцы метаду моцнай сувязі. \textbf{Прадмет даследавання} --- энергетычны спектр, адпаведныя ўласныя станы дыскрэтнага гамільтаніяна і залежнасць \(E_0\) і \(dE_1\) ад памеру і формы мяжы. \textbf{Мэта работы} --- мадэляванне энергетычнага спектра такіх кропак і праверка фізічна матываванай сурагатнай апраксімацыі; хвалевыя функцыі разглядаюцца як уласныя станы спектральнай задачы.

\textbf{Метады даследавання:} разлік спектра і ўласных станаў у Kwant, прамавугольная кантрольная задача, праверкі \(E_{\mathrm{kin}}=E_0+4\), \(1/a^2\) і Бесселя, Ridge-рэгрэсія, малы мнагаслойны перцэптрон, схемы LOAO/LOARO і аналіз астаткаў.

\textbf{Асноўныя вынікі.} Разлічаны 140 геаметрый: \(n=\{1.2,2.0,3.0,4.0\}\), \(a=\{24,27,30,33,36\}\), \(r_{\mathrm{AR}}\in[0.67,1.0]\). Памылка бесселевай ацэнкі складае каля \(2.03\%\), максімальнае адхіленне праверкі маштабавання пры фіксаванай форме --- \(2.12\%\). МП з фізічна матываванымі прыкметамі паспяховы ў 10 з 16 ячэек: 3 з 8 для LOAO і 7 з 8 для LOARO. Зададзены крытэрый не выкананы, таму асноўнай мадэллю застаецца Ridge. Адносная MAE Ridge складае 2.01--2.60\% ад \(E_{\mathrm{kin}}\) для \(E_0\) і 0.78--1.15\% ад \(dE_1\). Глабальнае тлумачэнне астаткаў простымі дыягностыкамі мяжы не пацверджана; асобны сігнал знойдзены толькі для \(n=1.2\) (\(\rho\approx -0.416\), \(p\approx0.013\)). \textbf{Абмежаванні.} Вынікі адносяцца толькі да даследаванай сеткі параметраў і фіксаваных класаў \(n\); сурагатная мадэль не замяняе прамы разлік Kwant.
\par
